%
% 23-potenzreihen.tex -- Potenzreihen
%
% (c) 2026 Prof Dr Andreas Müller
%
\subsection{Potenzreihen
\label{buch:komplex:funktionen:subsection:potenzreihen}}
Da in $\mathbb{C}$ der Begriff der Konvergenz wohldefiniert
ist, kann man eine komplexe Funktion auch durch eine konvergente
Potenzreihe definieren, wie man sie auch in der reellen Analysis
verwendet.

\begin{definition}[Potenzreihe]
\index{Potenzreihe}%
Eine \emph{Potenzreihe im Punkt $z_0\in\mathbb{C}$} ist eine Reihe der
Form
\begin{equation}
f(z)
=
\sum_{k=0}^\infty a_k(z-z_0)^k
\label{buch:komplex:funktionen:eqn:potenzreihe}
\end{equation}
mit $a_k\in\mathbb{C}$ für alle $k\in\mathbb{N}$.
\end{definition}

Der Begriff der Konvergenz einer Potenzreihe wird wie in der
reellen Analysis über die Konvergenz der Partialsummenfolge
definiert.

\begin{definition}[Partialsumme, Konvergenz]
\index{Partialsumme}%
\index{Partialsummenfolge}%
Die \emph{Partialsummenfolge} einer Potenzreihe ist die Folge
\[
s_n(z)
=
\sum_{k=0}^n a_k(z-z_0)^k
\]
für alle $n\in\mathbb{N}$.
Eine Potenzreihe heisst im Punkt $x\in\mathbb{C}$ \emph{konvergent}, wenn
die Partialsummenfolge $s_n(z)$ konvergent ist.
\end{definition}

\begin{beispiel}
\label{buch:komplex:funktionen:bsp:geometrisch}
Für $|z|<1$ ist die geometrische Reihe
\index{geometrische Reihe}%
\index{Reihe!geometrisch}%
\[
f(z)
=
1+z+z^2+z^3+z^4+\dots
\]
ist eine konvergente Potenzreihe im Punkt $0$.
Für reelle Werte ist bereits die Formel
\[
f(z)
=
\frac{1}{1-z}
\]
der Funktion $f(z)$ bekannt.
Um zu zeigen, dass dies auch für komplexe Argumente stimmt,
bestimmen wir die Partialsummen
\[
s_n(z)
=
1 + z + z^2 + z^3 + \dots + z^n.
\]
Sie erfüllen
\[
zs_n(z) - s_n(z)
=
z^{n+1}-1
\qquad\Rightarrow\qquad
s_n(z)
=
\frac{z^{n+1}-1}{z-1}.
\]
Die Differenz zweier solcher Folgenglieder ist
\[
s_n(z) - s_m(z)
=
\frac{z^{n+1}-1}{z-1}
-
\frac{z^{m+1}-1}{z-1}
=
\frac{z^{n+1}-z^{m+1}}{z-1}.
\]
Wir dürfen ohne Beschränkung der Allgemeinheit annehmen, dass $n<m$ ist.
Damit können wir die Differenz abschätzen
\[
|s_n(z)-s_m(z)|
=
\frac{1}{|z-1|}
|z|^{n+1} |1+z^{m-n}|
<
\frac{2}{|z-1|}
|z|^{n+1}.
\]
Die rechte Seite kann beliebig klein gemacht werden, denn ist
$\varepsilon>0$, dann wählen wir $n$ so, dass
\[
\frac{2}{|z-1|} |z|^{n-1}
<
\varepsilon
\quad\Rightarrow\quad
(n-1)\log |z|
<
\log\biggl(
\frac{\varepsilon}2 |z-1|
\biggr)
\quad\Rightarrow\quad
n
>
1
+
\frac{
\log \biggl(\frac{\varepsilon}{2}|z-1|\biggr)
}{\log|z|}
=
N(\varepsilon).
\]
Somit ist $s_n(z)$ ist eine Cauchy-Folge und damit ist die Reihe konvergent.
\end{beispiel}

In der reellen Analysis werden Kriterien für die Konvergenz von
Potenzreihen entwickelt.
Die dabei verwendeten Methoden bleiben für komplexe Argumente gültig.
Die Konvergenz wird im Allgemeinen auch von $z$ abhängen.
Im Beispiel~\ref{buch:komplex:funktionen:bsp:geometrisch} der
geometrischen Reihe ist die Konvergenz nur für $|z|<1$ gegeben.
Der maximale Betrag des Argumentes $z-z_0$, für den eine Potenzreihe
konvergiert, ist als Konvergenzradius bekannt.

\begin{definition}[Konvergenzradius]
\index{Konvergenzradius}%
Der \emph{Konvergenzradius} der Potenzreihe
\eqref{buch:komplex:funktionen:eqn:potenzreihe}
\[
\varrho
=
\sup
\biggl\{
|z-z_0|
\;
\bigg|
\;
\text{Es gibt ein $z\in\mathbb{C}$ derart, dass $\displaystyle
\sum_{n=0}^\infty a_n(z-z_0)^n$
konvergiert}
\biggr\}.
\]
Falls sich für jede reelle Zahl ein $z$ finden lässt, so dass die
Potenzreihe konvergiert, dann sagt man, dass der Konvergenzradius
unendlich ist.
\end{definition}

Die bekannten Kriterien der reellen Analysis für die Konvergenz
einer Reihe können dazu verwendet werden, den Konvergenzradius
zu berechnen.
Das Wurzelkriterium besagt zum Beispiel, dass eine Reihe mit
Summanden $a_n$ absolut konvergiert, wenn 
\[
\limsup_{n\to\infty} \!\sqrt[n]{|a_n|} < 1.
\]
Daraus ergibt sich die Cauchy-Hadamard-Formel
\begin{equation}
\varrho
=
\frac{1}{\displaystyle\limsup_{n\to\infty}\!\sqrt[n]{|a_n|}}.
\label{buch:komplex:funktionen:eqn:cauchy-hadamard}
\end{equation}
für den Konvergenzradius.
Im Beispiel~\ref{buch:komplex:funktionen:bsp:geometrisch}
ist $a_n=1$ für alle $n$, es folgt daher, dass der
Konvergenzradius $\varrho=1$ ist.

Falls alle bis auf endlich viele Koeffizienten einer Potenzreihen $\ne 0$
sind, kann auch das Quotientenkriterium angewendet werden.
Falls der Grenzwert
\[
r
=
\lim_{n\to\infty}
\biggl|
\frac{a_n}{a_{n+1}}
\biggr|
\]
existiert, ist $\varrho=r$ der Konvergenzradius.
Auch diese Kriterium liefert im 
Beispiel~\ref{buch:komplex:funktionen:bsp:geometrisch}
wieder den Konvergenzradius 1.

\begin{beispiel}
\label{buch:komplex:funktionen:bsp:exponentialfunktion}
Die Exponentialreihe~\eqref{buch:holomorph:holomorph:eqn:exponentialfunktion},
\index{Exponentialreihe}%
die in Kapitel~\ref{chapter:holomorph} genauer untersucht wird,
hat die Form
\[
\sum_{k=0}^\infty
\frac{1}{k!} z^k
\qquad\Rightarrow\qquad
a_k = \frac{1}{k!}.
\]
Zur Bestimmung des Konvergenzradius mit
dem Quotientenkriterium muss der Grenzwert
\[
\lim_{n\to\infty}
\biggl|
\frac{a_n}{a_{n+1}}
\biggr|
=
\lim_{n\to\infty}
\biggl|
\frac{(n+1)!}{n!}
\biggr|
=
\lim_{n\to\infty} n
=
\infty
\]
berechnet werden,
der Konvergenzradius ist also unendlich.

Dies zeigt sich auch in der Cauchy-Hadamard-Formel:
\index{Cauchy-Hadamard-Formel}%
\[
\limsup_{n\to\infty}
\!\sqrt[n]{|a_n|}
=
\limsup_{n\to\infty}
\!\sqrt[n]{\frac1{n!}}
\approx
\limsup_{n\to\infty}
\!\sqrt[n]{\frac{1}{\!\sqrt{2\pi n}}\biggl(\frac{e}{n}\biggr)^n}
=
\limsup_{n\to\infty}
\frac{1}{\!\sqrt[2n]{2\pi n}}
\frac{e}{n}
=
0,
\]
wobei wir die Stirling-Approximation $n!\approx\!\sqrt{2\pi n}(n/e)^n$
\index{Stirling-Approximation}%
\index{n!@$n!$}%
verwendet haben.
Auch mit diesem Kriterium ergibt sich also, dass die Reihe überall
in $\mathbb{C}$ konvergiert.
\end{beispiel}

